%!TEX root = ED_Grafo.tex

\section{Implementação}
% \begin{frame}[label=notas]
% 	\framesubtitle{Matriz de Adjacência}
	
% 	\begin{itemize}
% 		\item Cada posição na matriz é um arco:
% 			\begin{itemize}
% 				\item linha indica o nó de saída do arco
% 				\item coluna indica o nó de chegada
% 			\end{itemize}
% 		\pause
% 		\item Valor da posição indica:
% 			\begin{itemize}
% 				\item a existência ou não do arco
% 				\item o peso associado ao arco
% 			\end{itemize}
% 		\pause
% 		\item Deve ser utilizada para grafos densos, onde o número de arcos $|A|$ 
% 		é próximo do número de vértices ao quadrado $|V|^2$
% 		\subitem{Complexidade de espaço $O(V^2)$}
% 		\item O tempo necessário para acessar um elemento é  constante.
% 		\item Ler ou examinar  a matriz tem complexidade de 
% tempo $O(V^2)$
% 	\end{itemize}
% \end{frame}

\begin{frame}[label=notas]
	\framesubtitle{Matriz de Adjacência}

	\begin{columns}[T]
     	\column{.5\textwidth}\centering
     		\includegraphics[width=.8\textwidth]{graph_matrix1}
     	\pause
     	\column{.5\textwidth}\centering
     		\includegraphics[width=.8\textwidth]{graph_matrix2}
	\end{columns}
\end{frame}

% \begin{frame}[fragile,label=notas]
%	\framesubtitle{Lista de Adjacência}
 	
%	\begin{itemize}
%		\item Os vértices de uma lista de adjacência são em geral armazenados em 
%		uma ordem arbitrária.
%		\item Possui uma complexidade de espaço $O(|V|+|A|)$
%		\item Indicada para grafos esparsos, onde $|A|$ é muito menor do que $|V|^2$ .
%		\item é compacta e usualmente utilizada na maioria das aplicações.
%		\item Tempo $O(|V|)$ para determinar se existe uma aresta entre 
%		o vértice $i$ e o vértice $j$, pois podem existir $O(|V|)$vértices na 
%		lista de adjacentes do vértice $i$.
%	\end{itemize}
%	\begin{center}
%	\begin{minipage}{.35\textwidth}
%	\begin{lstlisting}[title=Vértice]
% typedef struct v_node_t {
%     char cidade[20];
%     a_node_t *aresta;
%     struct v_node_t *prox;
% } v_node_t;
%	\end{lstlisting}
%	\end{minipage}
%	\hspace{1cm}
%	\begin{minipage}{.4\textwidth}
%	\begin{lstlisting}[title=Aresta]
% typedef struct a_node_t {
%     double peso;
%     struct v_node_t *vertice;
%     struct a_node_t *prox;
% } a_node_t;
%	\end{lstlisting}
%	\end{minipage}
%	\end{center}
% \end{frame}

\makeatletter%
\newcommand{\vertice}[2][\@empty]{%
	\ifx#1\@empty%
		\node[block, fill=greenUnB, minimum size=2em] (#2) {\textcolor{white}{\large\textbf{#2}}};
	\else%
		\node[block, fill=greenUnB, minimum size=2em, below of=#1] (#2) {\textcolor{white}{\large\textbf{#2}}} edge[line, greenUnB, <-] (#1);
	\fi%
}%
\makeatother%
\newcommand{\aresta}[3]{%
	\node[block, blueUnB, minimum size=2em, right of=#1] (#1#2) {$#3 \overrightarrow{#2}$} edge[line,blueUnB,<-] (#1);
}%
\begin{frame}[label=notas]
	\framesubtitle{Lista de Adjacência}
	
	\begin{tikzpicture}[scale=1.5]
		\drawABCDnodes
		\drawABCDdirectedEdges
		\drawABCDweights
	\end{tikzpicture}
	\hfill\pause
	\begin{tikzpicture}[node distance=15mm]
		\vertice{A}
		\vertice[A]{B}
		\vertice[B]{C}
		\vertice[C]{D}
		
		\pause
		\aresta{A}{B}{0.3}
		\draw<3>[->, dashed] (AB)..controls +(south east:0.5) and +(north east:1) .. (B);
		\pause
		\aresta{B}{A}{0.4}
		\aresta{BA}{C}{0.2}
		\aresta{BAC}{D}{0.7}
		\draw<4>[->, dashed] (BA)..controls +(north east:0.5) and +(south east:1) .. (A);
		\draw<4>[->, dashed] (BAC)..controls +(south east:0.5) and +(north east:1) .. (C);
		\draw<4>[->, dashed] (BACD)..controls +(south east:0.5) and +(north east:1) .. (D);
		
		\pause
		\aresta{C}{C}{0.5}
		\aresta{CC}{D}{0.1}
		\aresta{D}{A}{0.8}
		\aresta{DA}{C}{1}
	\end{tikzpicture}
\end{frame}